\chapter{Khi Lớp Quá Ồn}
\chapterintro{Cắt ồn mà không làm mặt lớp đóng lại}{Vài câu đúng nhịp sẽ hiệu quả hơn rất nhiều lời quát nối nhau.}

\section{Lớp mình đông tiếng quá, nhường cho ý nào quan trọng nhất thôi}
\begin{manthoai}{Màn thoại}
\textbf{Thầy/Cô:} Lớp mình đông tiếng quá, nhường cho ý nào quan trọng nhất thôi.\\
\textbf{Học sinh:} Dạ.
\end{manthoai}
\begin{dungkhinào}
Dùng khi lớp đang ồn vì nhiều em cùng muốn nói chứ chưa phải quậy phá.
\end{dungkhinào}
\begin{nhipthem}
Kèm ngay luật: từng bạn một câu, nói xong mới đến lượt khác.
\end{nhipthem}
\ngancach

\section{Ai nói tiếp thì phải nối ý bạn trước}
\begin{manthoai}{Màn thoại}
\textbf{Thầy/Cô:} Ai nói tiếp thì phải nối ý bạn trước, không mở chương mới nhé.\\
\textbf{Học sinh:} Dạ.
\end{manthoai}
\begin{dungkhinào}
Hợp khi lớp ồn vì tranh nhau nêu ý nhưng thiếu nghe nhau.
\end{dungkhinào}
\begin{nhipthem}
Câu này biến ồn thành nghe. Nếu em nào không nối được, mời em chờ vòng sau.
\end{nhipthem}
\ngancach

\section{Giờ mình thử nói ít mà trúng}
\begin{manthoai}{Màn thoại}
\textbf{Thầy/Cô:} Giờ mình thử nói ít mà trúng, xem lớp có làm được không.\\
\textbf{Học sinh:} Có ạ.
\end{manthoai}
\begin{dungkhinào}
Hợp với lớp thích thi thố. Biến việc giảm ồn thành một thử thách vui sẽ hiệu quả hơn mệnh lệnh khô.
\end{dungkhinào}
\begin{nhipthem}
Nói xong thì chuyển ngay sang hoạt động thật ngắn để lớp có chỗ thử cách mới.
\end{nhipthem}
\ngancach

\section{Mình hạ âm lượng chứ không hạ ý kiến nhé}
\begin{manthoai}{Màn thoại}
\textbf{Thầy/Cô:} Mình hạ âm lượng chứ không hạ ý kiến nhé.\\
\textbf{Học sinh:} Dạ.
\end{manthoai}
\begin{dungkhinào}
Hợp khi lớp đang ồn vì tranh luận có thật, không phải vì phá rối.
\end{dungkhinào}
\begin{nhipthem}
Câu này giúp học sinh thấy mình không bị bịt miệng, chỉ đang được chỉnh cách nói.
\end{nhipthem}
\ngancach

\section{Lớp thử im đúng 5 giây để xem còn nghe được gì}
\begin{manthoai}{Màn thoại}
\textbf{Thầy/Cô:} Lớp thử im đúng 5 giây để xem còn nghe được gì.\\
\textbf{Học sinh:} Dạ.
\end{manthoai}
\begin{dungkhinào}
Rất hợp khi lớp ồn liên tục và cần một điểm ngắt ngắn nhưng không nặng nề.
\end{dungkhinào}
\begin{nhipthem}
Sau 5 giây đó, chốt ngay luật nói tiếp để lớp có đường quay trở lại hoạt động.
\end{nhipthem}
